\section{Conclusion}\label{sec:conclusion}
This review examined graph-based robotic world models for dynamic environments through three central questions: which representations are suitable, how they can be constructed and maintained, and how they can be exploited for downstream robotic tasks.
To support this analysis, we introduced a taxonomy organized around the temporal horizon of environmental dynamics and assessed factor graphs, scene graphs, and their hybrid combinations.
The surveyed literature shows that no single representation dominates across all forms and timescales of dynamics. 
Factor graphs provide a principled foundation for uncertainty-aware geometric estimation, whereas scene graphs enable semantic abstraction and relational reasoning.  Hybrid architectures can combine these complementary capabilities within a common world model, although the surveyed evidence does not establish their general superiority over alternative representations.

% \textcolor{red}{\sout{Construction and maintenance pipelines are increasingly mature for single-session operation under moderate dynamics. Nevertheless, scalable lifelong maintenance, robust open-vocabulary data association under severe dynamics, consistent uncertainty propagation across geometric and semantic layers, and distributed multi-robot construction remain unresolved.
% Although existing representations already support a broad range of downstream capabilities, including navigation, manipulation, prediction, and task planning, their reliability ultimately depends on maintaining consistency and quantifiable confidence as the environment evolves. 
% Overall, tighter integration among probabilistic factor graphs, semantically rich scene graphs, and emerging dense geometric representations constitutes a central research direction toward robust situational awareness}}
The open challenges analyzed in Section X reveal a dependency structure rather than a set of equally independent research directions.
At the representation and estimation level, reliable uncertainty handling, robust data association, and principled entity-state management are required to maintain a consistent dynamic world model.
Extending this consistency across longer temporal horizons and multiple robots introduces a second set of challenges related to lifelong maintenance, scalability, and distributed construction.
Downstream capabilities such as uncertainty-aware planning, active perception, and task reasoning build on these maintained representations, while introducing additional task-level challenges of their own.
Task-conditioned evaluation cuts across all three levels and is necessary to determine whether improvements in representation and maintenance translate into measurable gains in robotic performance.
Overall, the surveyed evidence indicates that progress toward robust long-term dynamic world models depends on strengthening these representation and maintenance foundations while developing evaluation protocols that expose their impact on downstream autonomy in dynamic environments.

% The open challenges above allow the three research questions presented in Section \ref{sec:introduction} to be summarized. 
% Regarding \textit{which representations} are suited to dynamic environments, the taxonomy shows that no single graph stack dominates; instead, key limitations arise at the interfaces between paradigms (notably uncertainty transfer across the FG-SG boundary and incomplete coverage across dynamics timescales). 
% Regarding \textit{how representations can be constructed and maintained}, the pipeline is increasingly well understood for single sessions with moderate dynamics, but scalable lifelong maintenance, robust open-vocabulary association under severe dynamics, and multi-robot shared construction remain unresolved. 
% Regarding \textit{how representations can be exploited}, downstream capabilities are already broad, yet their reliability depends on representations that remain consistent under change and can quantify confidence across geometric and semantic layers. 
% In this context, tighter integration between probabilistic factor graphs, semantically rich scene graphs, and emerging dense geometric substrates, in our view, one of the central challenges, and a promising direction, for long-term robotic situational awareness in dynamic environments.